\chapter{Sản phẩm: Dashboard hỗ trợ điều phối}

\section{Vai trò của giao diện trong DSS}

Một DSS chỉ có model nhưng không có giao diện quyết định sẽ khó được sử dụng.
Dashboard là lớp tương tác giữa dữ liệu, mô hình và người điều phối. Nó không
chỉ hiển thị Accuracy; nó phải cho phép người dùng:

\begin{itemize}[leftmargin=*]
  \item quan sát workload và phân bố Priority;
  \item truy vết từ khuyến nghị về xác suất, category và conflict;
  \item lọc theo event, đội nhận và nhu cầu review;
  \item thử một trường hợp mới;
  \item xuất hàng đợi để tiếp tục xử lý.
\end{itemize}

Streamlit là frontend của prototype. Đây là framework Python tạo giao diện web
từ widget, biểu đồ và state; trình duyệt là phía hiển thị, còn mã Python chạy ở
server. Vì vậy dự án không cần thêm React hoặc một frontend riêng để đáp ứng
luồng sử dụng học thuật. Frontend tách rời chỉ cần khi có API độc lập, phân
quyền phức tạp, nhiều người dùng đồng thời hoặc yêu cầu giao diện cao hơn.

\section{Năm trang chức năng}

\begin{description}[leftmargin=3.8cm,style=nextline]
  \item[Overview] Hiển thị số case, Priority, Manual Review, routing team,
        category và phân bố Risk Score. Bộ lọc event giúp quan sát thay đổi
        workload giữa các thảm họa.
  \item[EDA] Trình bày thống kê train-only, phân bố nhãn, K-Means, Apriori,
        kiểm kê ảnh và bất đồng đa phương thức. Trang này giải thích căn cứ
        thiết kế, không dùng test để khám phá.
  \item[Model Evaluation] Hiển thị so sánh sáu model trên dev, tuning,
        baseline, fusion, robust test và metric từng lớp. Người đọc có thể thấy
        cả điểm mạnh lẫn lớp yếu.
  \item[Single-Case Demo] Nhận tweet và ảnh upload. Chế độ text-only cho kết quả
        nhanh; chế độ text+image chạy CLIP rồi toàn bộ fusion--DSS. Kết quả gồm
        xác suất, category, conflict, Risk, Priority, team và action.
  \item[Priority Queue] Cho phép lọc theo Priority, team, event và Manual Review;
        xem chi tiết từng hàng và tải CSV để bàn giao.
\end{description}

\begin{figure}[H]\centering
\begin{subfigure}{0.48\textwidth}
\includegraphics[width=\textwidth]{01_overview_safe.png}
\caption{Overview và workload.}
\end{subfigure}
\begin{subfigure}{0.48\textwidth}
\includegraphics[width=\textwidth]{02_eda_train_only.png}
\caption{EDA train-only.}
\end{subfigure}
\caption{Tổng quan vận hành và căn cứ dữ liệu.}
\end{figure}

\begin{figure}[H]\centering
\begin{subfigure}{0.48\textwidth}
\includegraphics[width=\textwidth]{03_model_fusion_safe.png}
\caption{Đánh giá mô hình.}
\end{subfigure}
\begin{subfigure}{0.48\textwidth}
\includegraphics[width=\textwidth]{05_priority_queue_safe.png}
\caption{Priority Queue.}
\end{subfigure}
\caption{Từ bằng chứng mô hình tới hàng đợi quyết định.}
\end{figure}

\section{Kiến trúc dữ liệu của dashboard}

Dashboard không huấn luyện lại model khi người dùng mở trang. Batch test đã
được xử lý thành \code{dashboard\_test\_predictions.csv}, gồm xác suất fusion,
category, Risk Score, Priority, routing và cờ review. Cách cache này:

\begin{itemize}[leftmargin=*]
  \item giảm thời gian tải trang;
  \item bảo đảm biểu đồ dùng đúng model đã khóa;
  \item tách đánh giá hàng loạt khỏi suy luận một trường hợp.
\end{itemize}

Single-Case Demo tải vectorizer và bốn classifier. Model ảnh được lazy-load:
chỉ khi người dùng upload ảnh thì PyTorch/Transformers và CLIP mới được nạp.
Nhờ vậy chế độ text-only khởi động nhanh và cloud không phải tải CLIP ở mọi
trang.

\section{Công nghệ và lý do lựa chọn}

\begin{table}[H]\centering\small
\caption{Vai trò của các công nghệ chính.}
\begin{tabularx}{\textwidth}{p{3.1cm}p{3.3cm}X}
\toprule
\textbf{Công nghệ} & \textbf{Định nghĩa/vai trò} & \textbf{Lý do sử dụng}\\
\midrule
Python & Ngôn ngữ thực thi pipeline & Hệ sinh thái dữ liệu/ML thống nhất, dễ tái lập.\\
pandas/NumPy & Bảng dữ liệu và đại số mảng & Join annotation, audit, cache và tính metric.\\
scikit-learn & Thư viện ML cổ điển & Cung cấp TF--IDF, sáu classifier, metric và API nhất quán.\\
PyTorch + Transformers & Runtime mô hình tiền huấn luyện & Nạp CLIP và chạy forward pass ảnh.\\
Streamlit & Framework web bằng Python & Tạo frontend trực tiếp trên cùng logic dự án.\\
Plotly/Matplotlib & Thư viện trực quan & Biểu đồ tương tác trên app và hình tĩnh cho báo cáo.\\
Jupyter & Notebook thực thi & Lưu code, output và hình EDA/model trong một tài liệu.\\
XeLaTeX & Hệ dàn trang báo cáo & Hỗ trợ tiếng Việt, công thức, bảng và PDF nhất quán.\\
\bottomrule
\end{tabularx}
\end{table}

Việc dùng một stack Python giảm nguy cơ giao diện tính lại metric khác backend.
Đánh đổi là Streamlit không tối ưu cho ứng dụng giao dịch thời gian thực hoặc
UI tùy biến sâu.

\section{Chạy local và triển khai Streamlit Community Cloud}

\subsection{Chạy local đầy đủ}

\begin{verbatim}
pip install -r requirements.txt
streamlit run app/streamlit_app.py
\end{verbatim}

Nếu cần tái tạo toàn bộ dữ liệu, embedding và mô hình:

\begin{verbatim}
python -m scripts.run_all
\end{verbatim}

Lệnh tái tạo đầy đủ cần corpus ảnh CrisisMMD khoảng 1,8 GB và thời gian chạy
CLIP. Khi model/cache đã có, không cần chạy lại pipeline trước khi mở dashboard.

\subsection{Bundle cloud}

Repository cloud chứa:

\begin{itemize}[leftmargin=*]
  \item processed train/test cần cho EDA và dashboard;
  \item vectorizer và bốn classifier đã khóa;
  \item cấu hình fusion và các bảng metric;
  \item mã Streamlit cùng requirements CPU.
\end{itemize}

Raw image corpus và ma trận embedding lớn không được đưa lên Git. Trang EDA đọc
processed CSV. Single-Case Demo nhận ảnh upload; lần suy luận ảnh đầu tiên tải
CLIP nên có cold start, còn text-only dùng ngay các model nhỏ.

Cấu hình deploy:

\begin{center}
\begin{tabular}{ll}
\toprule
Repository & \code{dafiiscoding/ml-dss}\\
Branch & \code{main}\\
Main file & \code{app/streamlit\_app.py}\\
Python & 3.12\\
Secrets & Không yêu cầu\\
\bottomrule
\end{tabular}
\end{center}

Các bước trên được mô tả chi tiết trong \code{docs/DEPLOY\_STREAMLIT.md}.

\section{Kiểm thử giao diện và giới hạn triển khai}

Sáu entrypoint Streamlit gồm trang chính và năm trang chức năng được chạy bằng
AppTest trong môi trường cloud Python 3.12. GitHub Actions kiểm tra bundle,
dependency và text-only inference sau mỗi lần push. Unit test kiểm tra logic
dữ liệu, duplicate mask, tuning, robustness và DSS; integration test chạy một
ảnh thật qua CLIP.

Những kiểm thử này chứng minh app khởi động và luồng chính không phát sinh
exception. Chúng chưa chứng minh khả năng chịu tải nhiều người dùng, bảo mật
upload, phân quyền, lưu audit log hoặc SLA. Nếu triển khai thực tế cần thêm API,
authentication, giới hạn file, quét nội dung, logging, giám sát latency và chính
sách bảo vệ dữ liệu.

\begin{insight}[Phạm vi sản phẩm]
Dashboard là một prototype DSS có frontend hoàn chỉnh cho học tập và trình diễn:
người dùng có thể xem bằng chứng, thử ca mới và thao tác hàng đợi. Nó chưa phải
hệ thống chỉ huy production và không được dùng để tự động điều động nguồn lực.
\end{insight}
